% ============================================================
% PARTE MARIANO — tasa de descuento y saltos en Vt;
% más adelante, papers posteriores
% (implementadas las 4 primeras slides finales del borrador de Mariano)
% Fórmulas cotejadas con McDonald-Siegel (1986), ecs. (8),(10),(12),(13),(14).
% ============================================================
\divisoria{Mariano Sanchez}% TEMPORAL: borrar en la versión final

% M1 — cálculo de la tasa de descuento (Itô, ec. 8)
\begin{frame}{Cálculo de la tasa de descuento}
  Hasta acá se dio por sentada la tasa \az{$\mu$} a la que se descuentan los
  flujos. Se demuestra que la tasa esperada de equilibrio de la oportunidad
  debe ser un \concepto{promedio ponderado} de las tasas de activos con el
  mismo riesgo que \rj{$V$} y \az{$F$}.
  \par\vspace{2mm}
  Aplicando el lema de Itô al valor de la opción \rj{$X$}:
  \begin{ecmed}
    \rj{\frac{dX}{X}} = \az{\bigl[\varepsilon\alpha_v + (1-\varepsilon)\alpha_f
      + \tfrac12\varepsilon(\varepsilon-1)\sigma^2\bigr]dt}
      + \underbrace{\az{\varepsilon\sigma_v\,dz_v
      + (1-\varepsilon)\sigma_f\,dz_f}}_{\text{componente imprevisto}}
  \end{ecmed}
  con \az{$\sigma^2 = \sigma_v^2 + \sigma_f^2 - 2\rho_{vf}\sigma_v\sigma_f$}.
  El retorno de la opción es un promedio ponderado (\az{$\varepsilon$},
  \az{$1-\varepsilon$}) de los de \rj{$V$} y \az{$F$}.
\end{frame}

% M2 — aplicación del CAPM -> mu
\begin{frame}{Aplicación del CAPM al riesgo}
  El único riesgo de la opción es su componente imprevisto. Bajo el CAPM, un
  riesgo \az{$\sigma_j\,dz_j$} exige una tasa \az{$\hat{\alpha}_j$}, y un
  riesgo \az{$\nu\sigma_j\,dz_j$} exige \az{$r + \nu(\hat{\alpha}_j - r)$}.
  \par\vspace{3mm}
  Aplicándolo a los pesos \az{$\varepsilon$} y \az{$1-\varepsilon$}, la tasa
  de rendimiento esperada de equilibrio es
  \begin{ec}
    \rj{\mu} = \az{r + \varepsilon(\hat{\alpha}_v - r)
      + (1-\varepsilon)(\hat{\alpha}_f - r)
      = \varepsilon\hat{\alpha}_v + (1-\varepsilon)\hat{\alpha}_f}
  \end{ec}
\end{frame}

% M3 — epsilon en términos de delta
\begin{frame}{La solución para $\varepsilon$ en términos de $\delta$}
  Igualando el rendimiento requerido con el efectivo, la solución cuadrática
  para \rj{$\varepsilon$} queda en términos de los costos de oportunidad
  \az{$\delta_v = \hat{\alpha}_v - \alpha_v$}, \az{$\delta_f = \hat{\alpha}_f
  - \alpha_f$}:
  \begin{ecmed}
    \rj{\varepsilon} = \az{\sqrt{\left(\frac{\delta_f-\delta_v}{\sigma^2}
      - \frac12\right)^{\!2} + \frac{2\delta_f}{\sigma^2}}
      + \left(\frac12 - \frac{\delta_f-\delta_v}{\sigma^2}\right)}
  \end{ecmed}
  \begin{itemize}
    \item Los drift \az{$\alpha$} \resultado{no importan por sí solos}: lo que
          entra en la solución son los \concepto{costos de oportunidad}
          \az{$\delta_v$}, \az{$\delta_f$}.
    \item Dos proyectos con riesgo sistemático opuesto pueden valer lo mismo,
          según los estados en que generen rendimiento.
  \end{itemize}
\end{frame}

% M4 — análisis de los parámetros delta
\begin{frame}{Los parámetros $\delta_v$ y $\delta_f$}
  \concepto{Costo de mantener la opción} (\az{$\delta_v$}):
  \begin{itemize}
    \item Parte del rendimiento exigido de \rj{$V$} a la que se renuncia al
          percibir solo aumentos de precio.
    \item Si \az{$\delta_v$} sube: la inversión óptima se adelanta (menor
          umbral \rj{$V$}$/$\az{$F$}) y la opción \resultado{vale menos}.
    \item Más riesgo sistemático de \rj{$V$} (con \az{$\alpha_v$} constante)
          eleva \az{$\delta_v$}.
  \end{itemize}
  \par\vspace{2mm}
  \concepto{Beneficio de postergar} (\az{$\delta_f$}):
  \begin{itemize}
    \item Rendimiento de \az{$F$} al que se renuncia al obtener solo aumentos
          de precio (costo aplazado).
    \item Si \az{$\delta_f$} sube: mayor beneficio por esperar, la opción
          \resultado{vale más}. Más riesgo sistemático de \az{$F$} eleva
          \az{$\delta_f$}.
  \end{itemize}
\end{frame}

% M5 — caso sin y con incertidumbre (ec. 13, 14, 14')
\begin{frame}{La regla de inversión: sin y con incertidumbre}
  \concepto{Sin incertidumbre} (\az{$\sigma_v^2 = \sigma_f^2 = 0$}):
  \begin{ecmed}
    \rj{\hat{X}} = \az{(\hat{C} - 1)\,F_0
      \left(\frac{V_0/F_0}{\hat{C}}\right)^{\!\delta_f/(\delta_f-\delta_v)}},
    \quad \az{\hat{C} = \delta_f/\delta_v}, \quad
    \text{\small invertir si } \az{V_0\delta_v \geq F_0\delta_f}
  \end{ecmed}
  \vspace{-1mm}
  {\small\noindent Válido solo si \az{$\delta_v < \delta_f$}.}
  \par\vspace{1mm}
  \concepto{Con incertidumbre} (\az{$\sigma^2 > 0$}): la regla incorpora un
  término extra \az{$h$}, el valor de esperar:
  \begin{ecmed}
    \az{V_0\,\delta_v \geq F_0\,\delta_f + h(V_0, F_0)}, \qquad \az{h > 0}
  \end{ecmed}
  Ese término \az{$h$} permite \resultado{captar ganancias y evitar pérdidas}
  --- es lo que agrega la opción de esperar.
\end{frame}
